% Transcription — fonds Alexandre Grothendieck, shelfmark 161-1, batch 1 (pages 1-19).
% Established from the facsimile archives/batches/161-1/batch-01.pdf (a mirror of
% https://grothendieck.umontpellier.fr/161-1.pdf), per the skill
% transcribe-grothendieck. Pages 2, 8 and 14 are blank versos — absent.
% Page 16 is a reused verso, pencil notes upside down on an entirely different
% subject (p-adic periods) — absent, as it should be. Page 18 is a blank verso
% carrying three scribbled arrows with no labels — absent.
% Pass: Fable 5 (claude-fable-5), 2026-08-08 — first pass, unchecked against the pages by a human.
\documentclass[11pt,a4paper]{article}
% Le préambule commun à toutes les transcriptions du fonds Grothendieck.
%
% Il tient en une idée : **l'incertitude se compose, elle ne se résout pas.**
% Une transcription qui lisse un mot douteux a détruit la seule chose pour
% laquelle on la faisait — un lecteur doit pouvoir savoir, à chaque endroit,
% ce qui était sur le feuillet et ce qui a été ajouté. Les six macros qui
% suivent sont donc l'appareil critique, et non de la décoration.
%
% Les mêmes commandes sont reconnues par `scripts/render.mjs`, qui produit la
% vue de lecture du site. Une macro ajoutée ici doit l'être aussi là-bas,
% faute de quoi le rendu échouera bruyamment — ce qui est voulu : un rendu
% silencieusement faux est pire qu'un rendu absent.

\ProvidesPackage{grothendieck}[2026/08/08 Transcriptions du fonds Grothendieck]

\usepackage{amsmath,amssymb,amsthm}
\usepackage{mathrsfs}
\usepackage{stmaryrd}
% Les diagrammes commutatifs sont partout dans ce fonds : c'est la langue
% même de la géométrie algébrique de Grothendieck, et une transcription qui
% les rendrait en prose ne transcrirait rien.
\usepackage{tikz-cd}
% Français par défaut — c'est la langue des feuillets — mais l'anglais est
% chargé aussi : la traduction et le résumé sont des documents anglais, et la
% typographie française (espaces avant les deux-points) y serait une faute.
% Ces documents basculent par \selectlanguage{english} en tête de corps.
\usepackage[english,french]{babel}
\usepackage{fontspec}
\usepackage{geometry}
\geometry{a4paper,margin=25mm,marginparwidth=28mm}
\usepackage{marginnote}
\usepackage{xcolor}
% ulem, not soul. `soul`'s \st and \ul are built on a character-by-character
% scan that XeLaTeX's font machinery breaks: the document fails with an
% undefined control sequence rather than degrading. ulem does the same two jobs
% and is Unicode-engine safe. `normalem` stops it hijacking \emph, which the
% transcriptions use for Grothendieck's own emphasis.
\usepackage[normalem]{ulem}
\usepackage{hyperref}
\hypersetup{colorlinks=true,linkcolor=black,urlcolor=black,pdfborder={0 0 0}}

% --- Métadonnées du lot -----------------------------------------------------
% Reprises telles quelles par le script de rendu, qui les affiche en tête de
% la vue de lecture. Elles fixent ce que le fichier prétend couvrir : sans
% elles, une transcription flotte sans point d'attache dans un fonds de seize
% mille feuillets.

\newcommand{\folder}[1]{\gdef\grothFolder{#1}}
\newcommand{\batch}[1]{\gdef\grothBatch{#1}}
\newcommand{\leaves}[2]{\gdef\grothFirst{#1}\gdef\grothLast{#2}}
\newcommand{\foldertitle}[1]{\gdef\grothFoldertitle{#1}}
\gdef\grothFolder{?}\gdef\grothBatch{?}\gdef\grothFirst{?}\gdef\grothLast{?}\gdef\grothFoldertitle{}

% --- Le résumé --------------------------------------------------------------
%
% Il ouvre la lecture modernisée, sous le titre « Résumé », et il est composé
% en petit. Ce n'est pas un ornement : c'est le seul passage du document qui
% ne soit pas des mathématiques, et un lecteur doit pouvoir le reconnaître
% avant de l'avoir lu — puis le sauter s'il connaît déjà le sujet, ou s'y
% arrêter s'il ne le connaît pas. Le filet qui le suit marque la reprise du
% corps.

\newenvironment{resume}
  {\small\setlength{\parskip}{0.45em}}
  {\par\medskip\hrule\medskip}

% --- Mots-clés ---------------------------------------------------------------
%
% En anglais, en clôture du résumé de la lecture modernisée : le vocabulaire
% moderne sous lequel chercher ce que les feuillets construisent. Le manifeste
% du site (`npm run manifest`) les extrait pour étiqueter la cote entière —
% cette ligne est la source unique des tags, il n'y a pas de fichier à côté.
\newcommand{\keywords}[1]{\par\smallskip\noindent
  {\footnotesize\textsc{Keywords} --- \textit{#1}}\par}

% --- L'appareil critique ----------------------------------------------------

% Le numéro de feuillet, en marge. C'est l'ancre qui permet de régler un
% désaccord sur une formule en regardant *un* feuillet plutôt qu'une chemise
% de six cents. Le site s'en sert aussi pour tourner les pages du fac-similé
% quand on fait défiler la transcription.
\newcommand{\leaf}[1]{%
  \marginnote{\footnotesize\color{gray!70}\textsf{#1}}%
  \label{leaf:#1}\ignorespaces}

% Illisible : on ne devine pas. Un mot inventé qui se lit comme les autres est
% le pire résultat possible.
\newcommand{\ill}{\textcolor{red!60!black}{[\,\ldots\,]}}

% Lecture douteuse : le mot est donné, et signalé comme incertain.
\newcommand{\uncertain}[1]{\uline{#1}}

% Ajout de l'éditeur — une lettre manquante, un mot sous-entendu.
\newcommand{\add}[1]{\textcolor{blue!50!black}{[#1]}}

% Biffé par Grothendieck lui-même. Ce qu'il a rayé fait partie du document :
% une rature dit souvent où la pensée a changé de direction.
\newcommand{\struck}[1]{\sout{#1}}

% Note du transcripteur, sur l'état matériel du feuillet.
\newcommand{\note}[1]{\par\smallskip{\small\color{gray!60!black}\textit{[#1]}}\par\smallskip}

% Note marginale de Grothendieck, distincte du corps du texte.
\newcommand{\marginal}[1]{\par\smallskip\noindent\hspace{1em}%
  {\small\color{gray!60!black}#1}\par\smallskip}

% --- Titre ------------------------------------------------------------------

\AtBeginDocument{%
  \begin{center}
    {\large\bfseries Fonds Alexandre Grothendieck --- cote n\textsuperscript{o}~\grothFolder}\\[2pt]
    {\itshape\grothFoldertitle}\\[4pt]
    {\small feuillets \grothFirst--\grothLast\ (lot \grothBatch)}\\[2pt]
    {\footnotesize\color{gray!60!black}Transcription établie d'après le fac-similé numérisé
     par l'Université de Montpellier.\\
     Les lectures incertaines sont soulignées, les illisibles notées [\,\ldots\,],
     les ajouts entre crochets.}
  \end{center}
  \vspace{1em}}

\endinput


\folder{161-1}
\batch{1}
\pages{1}{19}
\foldertitle{Catégories : notes manuscrites (s.d.)}

\begin{document}

\page{1}
Titre porté au crayon sur l'onglet de la chemise : « Catégories ».
\note{chemise du dossier}

\section*{Foncteurs adjoints (pages 3 à 7)}

\page{3}
En tête, à droite : « Foncteurs adjoints [demander Schanuel] ».

Soit
\[
E \underset{v}{\overset{u}{\rightleftarrows}} F,
\qquad i : \mathrm{id}_E \to vu, \quad j : uv \to \mathrm{id}_F,
\]
un couple de foncteurs adjoints.

a) Soit $E'$ une sous-catégorie pleine de $E$. Notons que pour
$X, Y \in \mathrm{Ob}\,E'$, l'application
\[
\mathrm{Hom}_{E'}(X,Y) \longrightarrow \mathrm{Hom}_{F}(uX, uY)
\quad (\simeq \mathrm{Hom}(X, vuY))
\]
\note{sous le premier membre, le manuscrit rappelle
$\mathrm{Hom}_{E'}(X,Y) = \mathrm{Hom}_{E}(X,Y)$}
se déduit de l'application $i(Y) : Y \to vu(Y)$ au moyen de
$\mathrm{Hom}(X, i(Y))$. Donc pour que ce soit bijectif
\uncertain{(resp. injectif)} pour $\forall X, Y \in \mathrm{Ob}\,E'$ — i.e.
$u|E'$ pl. fidèle (resp. fidèle) — il suffit que
$\forall Y \in \mathrm{Ob}\,E'$, $i(Y) : Y \to vu(Y)$ soit un isomorphisme
(resp. un \uncertain{mono}). C'est aussi nécessaire si on sait que
$vu(E') \subset \overline{E}'$. \note{la parenthèse qui suivait est raturée
au point d'être illisible ; la ligne porte plusieurs corrections
superposées} On trouve donc

(1) $u|E'$ pl. fidèle et $vu(E') \subset \overline{E}'$
$\iff$ $\forall Y \in \mathrm{Ob}\,E'$, $i(Y) : Y \xrightarrow{\ \sim\ } vu(Y)$.

\marginal{donner un exemple montrant que cette condition n'est pas
surabondante, avec $E'$ réduit à un élément}

b) Soit $F'$ une sous-catégorie pleine de $F$. Notons que pour
$P, Q \in \mathrm{Ob}\,F'$, l'application
\[
\mathrm{Hom}_{F'}(P,Q) \longrightarrow \mathrm{Hom}(vP, vQ)
\quad (\simeq \mathrm{Hom}(uvP, Q))
\]
\note{même rappel en dessous : $\mathrm{Hom}_{F'}(P,Q) = \mathrm{Hom}_F(P,Q)$}
se déduit des $j(P) : uv(P) \to P$ par $\mathrm{Hom}(j(P), Q)$. Donc pour que
$v|F'$ soit fidèle (resp. pl. fidèle), il suffit que
$\forall P \in \mathrm{Ob}\,F'$, $j(P)$ soit un épi. (resp. un isom). Donc on
a de même :

(2) $v|F'$ pl. fidèle et $uv(F') \subset \overline{F}'$
$\iff$ $\forall P \in \mathrm{Ob}\,F'$, $j(P) : uv(P) \xrightarrow{\ \sim\ } P$.

\note{suit une reprise de (2), encadrée d'un long trait et biffée, qui
précisait « strictement pleine $\subset E$ » et « strictement pleines »}

c) L'ensemble des $E'$ \add{strictement pleines} satisfaisant (1), et des
$F'$ \add{strictement pleines} satisfaisant (2), est en correspondance
biunivoque, en posant
\[
\alpha(E') = \text{image essentielle de } u|E' = \overline{u(E')},
\qquad
\beta(F') = \text{image essentielle de } v|F' = \overline{v(F')}.
\]
\note{« strictement pleines » est chaque fois une correction interlinéaire
de la main de l'auteur}

Si $E'$ et $F'$ sont des sous-catégories str. pleines de $E$ resp. $F$, les
conditions suivantes sont équivalentes :

\begin{enumerate}
\item $E'$ satisfait (1), $F'$ satisfait (2), et $E'$ et $F'$ se
correspondent comme ci-dessus ;
\item $u(E') \subset F'$, $v(F') \subset E'$, et $u'$ et $v'$ pl. fidèles ;
\item —— et $u'$ est une équivalence ;
\item —— et $v'$ est une équivalence ;
\item —— et $u'$ et $v'$ pl. fid.
\end{enumerate}

\note{en marge droite, une sixième condition esquissée reprend (1) et (2) :
« ……, et $\forall Y \in \mathrm{Ob}\,E'$, $i(Y)$… ;
$\forall P \in \mathrm{Ob}\,F'$, $j(P) : uv(P) \to$… » — la fin rognée par
le bord de la page}

\page{4}
d) \struck{Passons maintenant} Soit $E_0$ la sous-catégorie str. pleine de
$E$ formée des $Y \in \mathrm{Ob}\,E$ tels que
$i(Y) : Y \xrightarrow{\ \sim\ } vu(Y)$. C'est donc la plus grande des
sous-catégories pleines $E'$ de $E$ satisfaisant (1), et si $E'$ est une
sous-catégorie pleine de $E$, elle satisfait (1) ssi $E' \subset E_0$. On
définit de même la sous-catégorie str. pleine $F_0$ de $F$ formée des
$P \in \mathrm{Ob}\,F$ tels que $j(P) : uv(P) \xrightarrow{\ \sim\ } P$, une
sous-catégorie pleine $F'$ de $F$ satisfaisant (2) ssi $F' \subset F_0$.
Ceci posé, on a
\[
F_0 = \overline{u(E_0)}, \qquad E_0 = \overline{v(F_0)}
\]
(i.e. $F_0 =$ image ess. de $u|E_0$, $E_0 =$ image ess. de $v|F_0$) et $u$
et $v$ induisent sur $E_0$ et $F_0$ des équivalences inverses l'une de
l'autre.

e) Pour que l'on ait $v(F) \subset E_0$, il faut et il suffit que pour
$\forall P \in \mathrm{Ob}\,F$, on ait
\[
v(P) \xrightarrow{\ i(v(P))\ } vuv(P) \quad \text{(isomorphisme)}
\]
($v(j(P))$ sera alors l'inverse ??).

\marginal{En un mot, il faut et il suffit que $v$ se factorise en
$F \xrightarrow{v_1} E_1 \xrightarrow{\alpha} E$, avec $\alpha$ pl. fidèle,
$v_1$ passage à cat. des fractions. [NB nécessairement $v_1$ aura un adjoint
à g., \uncertain{sav.} $u_1 = u|E_1$, mais pas néc. $\alpha$]}

Dans ce cas, $E_0$ est l'image essentielle de $v$, et l'inclusion
$F_0 \subset F$ admet un adjoint \uncertain{à droite}, savoir
$u_0 \circ (v \text{ astreint à } E_0)$, noté $v'$ :
\note{la phrase est surchargée de ratures ; la lecture
« $F_0 \subset F$ … à droite » est celle qu'impose le calcul entre crochets
ci-dessous. Une première rédaction du calcul, barrée de grands traits
obliques, occupe le milieu de la page}

[En effet, on a, pour $Q_0 \in \mathrm{Ob}\,F_0$, $P \in \mathrm{Ob}\,F$ :
\[
\mathrm{Hom}_F(Q_0, P) \simeq \mathrm{Hom}_F(u_0 v_0 Q_0, P)
\quad \text{(car } u_0 v_0 \simeq \mathrm{id}\text{)}
\]
\[
\simeq \mathrm{Hom}_E(v_0 Q_0, vP) = \mathrm{Hom}_{E_0}(v_0 Q_0, v'P)
\simeq \mathrm{Hom}_{F_0}(u_0 v_0 Q_0, u_0 v' P)
\quad \text{($u_0$ pl. fid.)}
\]
\[
\simeq \mathrm{Hom}_{F_0}(Q_0, u_0 v' P)
\quad \text{(} u_0 v_0 \simeq \mathrm{id} \text{).]}
\]

Croquis en marge gauche, redessinés :

\begin{tikzcd}[column sep=small, row sep=small, nodes={font=\scriptsize}]
E \arrow[r, "u"] \arrow[d] & F \arrow[d] \arrow[dl, "v'"'] \\
E_0 \arrow[r, bend left=18, "u_0"] & F_0 \arrow[l, bend left=18, "v_0"]
\end{tikzcd}

\begin{tikzcd}[column sep=small, row sep=small, nodes={font=\scriptsize}]
E \arrow[r, "u"] \arrow[d] \arrow[dr] & F \arrow[d] \arrow[dl] \\
E_0 & F_0
\end{tikzcd}

\note{sur le premier croquis, $u_0$ et $v_0$ portent chacun une marque
$\approx$ — les équivalences quasi-inverses ; le second est griffonné, les
quatre flèches descendantes s'y croisent}

\note{suit, entre crochets et presque entièrement biffé, un passage qui pose
la question : si $\alpha v_1$ a un adjoint à g., avec $\alpha$ pl. fid. et
$v_1$ passage à cat. des fractions, « est-il vrai que $\alpha$ admette
adj. à g. ? » — le calcul qui l'accompagnait est barré}

Disons donc : pour que l'on ait $u(E) \subset F_0$, il faut et il suffit que
l'on ait pour $\forall X \in \mathrm{Ob}\,E$ :
\[
uvu(X) \xrightarrow{\ j(u(X))\ } u(X) \quad \text{(isomorphisme)},
\]
et alors $F_0$ est l'image ess. de $E$ par

\page{5}
$u$. Dans ce cas, l'inclusion $E_0 \subset E$ admet un adjoint à gauche,
savoir $v_0 \circ (u \text{ astreint à } F_0)$, noté $u'$ :
\[
\mathrm{Hom}_E(Y, X_0) \simeq \mathrm{Hom}_E(Y, v_0 u_0 X_0)
= \mathrm{Hom}_E(Y, v u_0 X_0)
\simeq \mathrm{Hom}_F(uY, u_0 X_0)
\]
\[
= \mathrm{Hom}_{F_0}(u'Y, u_0 X_0)
\simeq \mathrm{Hom}_{E_0}(v_0 u' Y, v_0 u_0 X_0)
= \mathrm{Hom}_{E_0}(v_0 u' Y, X_0). \qquad \text{cqfd}
\]
\note{sous les signes du calcul : « $v_0 u_0 \simeq \mathrm{id}$ »,
« adjonction », « $v_0$ pl. fid. », « $v_0 u_0 \simeq \mathrm{id}$ » ; sous
le premier membre : $Y \in \mathrm{Ob}\,E$, $X_0 \in \mathrm{Ob}\,E_0$}

f) Conditions équivalentes :

\begin{enumerate}
\item $\mathrm{Im}\,v \subset E_0$ et 1') $\mathrm{Im}\,u \subset F_0$ ;
\item $v \xrightarrow{\ \sim\ } vuv$ et 2') $uvu \xrightarrow{\ \sim\ } u$ ;
\item $\mathrm{Im}\,v \subset E_0$ \emph{et} l'inclusion $E_0 \subset E$
admet un adjoint à g. \note{3) s'ouvre sur une ligne entièrement biffée :
« …se factorise en un composé… » ; 3) et sa symétrique sont précédées d'un
crochet au crayon} — 3') $\mathrm{Im}\,u \subset F_0$ et l'inclusion
$F_0 \subset F$ admet un adjoint à dr. ;
\item $u$ se factorise en $E \xrightarrow{u_1} F_1 \xrightarrow{\beta} F$,
avec $\beta$ pleinement fidèle, $u_1$ \uncertain{passage à catégorie des
fractions}, et $\beta$ et $u_1$ ayant des adjoints à droite
\note{la seconde ligne de 4) est barrée d'un long trait qui traverse aussi
sa correction interlinéaire ; en marge gauche : « à \uncertain{revoir pour}
simplifier »} — 4') $v$ se factorise en
$F \xrightarrow{v_1} E_1 \xrightarrow{\alpha} E$, avec $\alpha$ pl. fidèle,
$v_1$ passage à cat. des fractions, [$v_1$ et $\alpha$ admettent un adj. à
g.].
\end{enumerate}

\struck{5)} De plus, la donnée d'un couple de foncteurs adjoints $(u,v)$
entre des catégories $E$, $F$, satisfaisant les conditions
\add{équivalentes} 1) à 4), équivaut à : la donnée d'une sous-cat. str.
pleine $E_0$ de $E$ telle que $\alpha : E_0 \subset E$ admette un adjoint à
g. $\alpha'$, d'une sous-cat. str. pleine $F_0$ de $F$ telle que
$\beta : F_0 \subset F$ admette un adjoint à droite $\beta'$, et d'une
équivalence entre $E_0$ et $F_0$
\[
E_0 \underset{v_0}{\overset{u_0}{\rightleftarrows}} F_0,
\]
en posant $u = \beta u_0 \alpha'$, $v = \alpha v_0 \beta'$, les données
d'adjonction étant évidentes.

\page{6}
En tête : $E \underset{v}{\overset{u}{\rightleftarrows}} F$.

a) $E$ catégorie ponctuelle, $v$ l'unique foncteur $F \to E$. Il a un
adjoint à g. ssi $F$ a un objet initial, et alors $u(P) = \emptyset_F$.
Alors $u$ est pl. fidèle.

b) $F$ catégorie ponctuelle, $u$ l'unique foncteur $E \to F$. Alors
l'adjoint à droite $v$ existe ssi $E$ a un objet final,
$v(P) = e_E$. Alors $v$ est pl. fidèle.

En marge gauche, le dispositif : $\alpha : E_1 \to E$ pl. fidèle, et son
adjoint à gauche $\alpha'$ :

\begin{tikzcd}[row sep=small]
E \arrow[d, bend right=25, "\alpha'"'] \\
E_1 \arrow[u, bend right=25, "\alpha"']
\end{tikzcd}

On suppose que $\alpha$ admet un adjoint à gauche $\alpha'$, d'où couple
$\alpha', \alpha$ ; \struck{qui fait} \uncertain{de sorte que} $\alpha'$ est
un passage à une catégorie des fractions. Considérons alors un autre
foncteur pl. fidèle $E_1 \xrightarrow{\beta} F$, et supposons que
$u = \beta\alpha'$ :

\begin{tikzcd}[column sep=small, row sep=small]
E \arrow[rr, "u"] \arrow[dr, "\alpha'"'] & & F \\
& E_1 \arrow[ur, "\beta"'] &
\end{tikzcd}

ait un adjoint à \struck{gauche} droite (cela exige de toutes façons,
\ill{} $\alpha'$, \uncertain{qu'existe un} adjoint à droite ; ici le fait
que $\alpha$ pl. fidèle équivaut à : $\alpha'$ passage à cat. des
fractions !). Est-il vrai alors que $\beta$ ait un adjoint à droite ? On
peut supposer que $\mathrm{Ob}\,F = \mathrm{Ob}\,E_1 \cup \{a\}$ ; il est
donné alors par deux foncteurs
\[
\varphi : E_1 \to (\mathrm{Ens}), \qquad
\psi : E_1^{\circ} \to (\mathrm{Ens}), \qquad \text{savoir}
\]
\[
\varphi(X) = \mathrm{Hom}(a, \beta X), \qquad
\psi(X) = \mathrm{Hom}(\beta X, a),
\]
et un morphisme bifonctoriel
\[
\varphi(X) \times \psi(Y) \longrightarrow \mathrm{Hom}(Y, X).
\]
\note{devant le premier membre, un début de formule est biffé ; une
parenthèse interlinéaire au-dessus de $\mathrm{Hom}(Y,X)$ reste illisible}

Dire que $\beta$ a un adjoint à droite signifie que le foncteur
contravariant $\psi$ est représentable. Prenant $\varphi(X) = \emptyset$
pour tout $X$, la donnée du morphisme bifonctoriel devient triviale.
Exprimons que $\beta\alpha'$ a un adjoint à droite, i.e. que le foncteur
sur $E$
\[
X \mapsto \mathrm{Hom}_F(\beta\alpha'(X), a) = \psi(\alpha'(X))
\]
est représentable. Donc la question devient simplement : un foncteur
$\psi : E_1^{\circ} \to (\mathrm{Ens})$ tel que $\psi \circ \alpha'$ soit
rep., est-il représentable ? \struck{Identifions} \uncertain{Identifiant les
$\psi$ aux foncteurs} sur $E^{\circ}$ qui transforment \uncertain{les
éléments} de $\Sigma$ en bij., on trouve : soit $\xi \in \mathrm{Ob}\,E$ tel
que le foncteur qu'il représente sur $E$, $X \mapsto \mathrm{Hom}_E(X,\xi)$,
transforme les $u \in \Sigma$ \struck{bijections} en bijections. Alors
est-il vrai que $\exists\, \xi_1 \in E_1$

\note{$\Sigma$ apparaît ici sans définition : c'est la classe des flèches
que $\alpha'$ inverse, comme l'explicite la page suivante}

\page{7}
tel qu'on ait un isom fonctoriel en $X$
\[
\mathrm{Hom}_E(X, \xi) \simeq \mathrm{Hom}_{E_1}(\alpha' X, \xi_1)
\quad (\simeq \mathrm{Hom}_E(X, \alpha\,\xi_1)),
\]
en d'autres termes a-t-on $\xi \in \mathrm{Ob}\,\overline{E}_1$ (image
essentielle de $\alpha$) ?

Or quelle est la relation entre $\Sigma$ et $E_1$ ?
\[
(u : X \to Y) \in \Sigma
\iff \alpha'(X) \simeq \alpha'(Y)
\iff \mathrm{Hom}_{E_1}(\alpha'(Y), Z_1) \simeq
     \mathrm{Hom}_{E_1}(\alpha'(X), Z_1)
\ \ \text{pour } \forall\, Z_1 \in \mathrm{Ob}\,E_1
\]
\note{au premier membre, \struck{« un isom »} est biffé et remplacé en
interligne par « $\alpha'(X) \simeq \alpha'(Y)$ », qui est la lecture donnée
ici}
\[
\iff \mathrm{Hom}_{E_1}(Y, \alpha Z_1) \simeq \mathrm{Hom}_{E_1}(X, \alpha Z_1)
\quad \text{pour } \forall\, Z_1 \in \mathrm{Ob}\,E_1.
\]
\note{au second membre, le manuscrit répète l'indice $E_1$ là où la lecture
$E$ s'impose — $\alpha Z_1$ est un objet de $E$}

Donc la question est si on récupère bien $\overline{E}_1$ comme
sous-catégorie str. pleine de $E$ comme la sous-catégorie str.
\struck{pleine} $\Phi(\Sigma) \subset E$ formée des $\xi$ tels que
$\mathrm{Hom}(Y, \xi) \simeq \mathrm{Hom}(X, \xi)$ pour
$\forall\, u : X \to Y$ dans $\Sigma$. Or c'est vrai sauf erreur…

\section*{Produit tensoriel et familles finies (pages 9 à 12)}

\page{9}
Soit $A$ une \struck{$\otimes$-}catégorie \struck{commutative}.
\note{une seconde ligne, entièrement biffée, lisait :
« \uncertain{une $\otimes$-catégorie AUC dont tous les objets inversibles} » ;
le mot « Définition » qui ouvrait le paragraphe suivant est biffé aussi}

Soit $\Phi(A)$ la catégorie des familles finies d'objets de $A$, un
morphisme de $(I, (L_i)_{i \in I})$ dans $(J, (M_j)_{j \in J})$ étant un
couple $(\tau, (\varphi_i)_{i \in I})$ d'une bijection
$\tau : I \xrightarrow{\ \sim\ } J$ et d'isom.
$\varphi_i : L_i \xrightarrow{\ \sim\ } M_{\tau(i)}$. On introduit sur
$\Phi(A)$ une $\otimes$-structure en posant
\[
(I, (L_i)_{i\in I}) \otimes (J, (M_j)_{j\in J})
= (I \sqcup J,\ (L_i;\ M_j)) \ \ldots
\]
une structure d'\struck{additivité} associativité, provenant de
l'associativité dans l'opération $I \sqcup J$, une structure unité par la
famille vide, étant aussi l'objet unité, une structure de commutativité par
la commutativité de $I \sqcup J$. (NB. pour ces définitions, la
$\otimes$-structure de $C$ n'a pas encore servi.)

\note{suit une phrase encombrée de ratures : « On \uncertain{va} trouver un
$\otimes$-foncteur AUC $\pi : \Phi(C) \to C$ tel que (la restriction de
$\pi$ aux… » — abandonnée au profit du corollaire ci-dessous ; en marge
droite, une note griffonnée : « $\Phi_n(C) = $ … »}

On a un foncteur canonique
\[
\Phi(A) \longrightarrow (\text{Ens finis, avec isom comme morphismes})
\]
(foncteur transportable …) ; on désigne par $\Phi_n(A)$ l'image inverse de
la sous-catégorie des ensembles de cardinal $n$ ($n$ entier $\geq 0$). On a
donc
\[
\Phi(A) = \coprod_n \Phi_n(A),
\qquad \Phi_0(A) = \{e\}, \quad \Phi_1(A) \simeq A.
\]

1) Soit $C$ une catégorie munie d'un produit tensoriel $\otimes$ AUC.
Considérons le foncteur « restriction à $\Phi_1(A)$ » :
\[
\underline{\mathrm{Hom}}^{\otimes \mathrm{AUC}}(\Phi(A), C)
\longrightarrow \underline{\mathrm{Hom}}(A, C).
\]

\textbf{Proposition.} C'est une équivalence de catégories.

\emph{Dém.} ….

\textbf{Corollaire.} Soit $C$ une $\otimes$-catégorie AUC. Il y a, à isom
unique près, un et un seul $\otimes$-foncteur AUC
\[
\pi : \Phi(C) \longrightarrow C, \qquad \pi = (\pi_n),
\]

\page{10}
muni d'un isom. rendant commutatif \note{le diagramme qui suivait —
$\Phi_1(C) \rightrightarrows C$, comparant $\pi_1$ au foncteur canonique —
est recouvert de gribouillis}

\textbf{Définition.} Si $(I, (L_i)_{i\in I})$ est une famille finie
d'objets de $C$, $\pi((L_i))$ est appelé le produit $\otimes$ de la famille
dans $C$.

Il y a une vérification : si on se donne une cat. et une opération
$\otimes$, alors la donnée d'une structure AUC sur celle-ci est équivalente
à la donnée d'un foncteur $\pi : \Phi(C) \to C$, \struck{et d'un isom.
$\pi(X \otimes Y) \simeq \pi(X) \otimes \pi(Y)$, en associant à la
structure d'associativité} qui sur $\Phi_1(C)$ soit « l'identité », et qui
rende commutatif de plus le diagramme

\begin{tikzcd}[column sep=small, row sep=small]
\Phi_2(C) \arrow[r, "\pi\vert\Phi_2"] & C \\
C \times C \arrow[u] \arrow[ur, "\otimes"'] &
\end{tikzcd}

\textbf{NB} Dans $\Phi(C)$ lui-même : si $(\xi_j)_{j\in J}$ est une famille
finie d'objets de $\Phi(C)$, $\xi_j = (I_j, (L_{ji})_{i \in I_j})$, alors
$\pi((\xi_j)_{j\in J})$ est
\[
\Bigl(\coprod_{j\in J} I_j,\ (L_{ji})_{j\in J,\ i\in I_j}\Bigr)
\]
(\uncertain{avec abus de notation}).

Revenons : la $\otimes$-cat $C$ AUC. Soit $I$ un ensemble à deux éléments,
$(L_i)_{i\in I}$ une famille d'objets de $C$ indexée par $I$. Il y a deux
ordres totaux sur $I$, et si les deux éléments de $I$ sont notés $a$, $b$,
ces deux façons de regarder le produit, suivant l'ordre total choisi, sont
$L_a \otimes L_b$ et $L_b \otimes L_a$ ; on passe de l'une à l'autre par
l'isom de symétrie. Cela offre \uncertain{deux données} si $L_a = L_b = L$
\note{la fin de la phrase, raturée, évoque $\varepsilon(L)$ et
l'automorphisme d'ordre 2 de $L \otimes L$ ; elle se poursuit à la page
12 — la page 11 porte une construction intercalée}

\page{11}
\textbf{Construction.} Soit $\Phi^{c}(C)$ la catégorie formée des
\struck{familles} couples $(\xi, \gamma)$, où $\xi \in \mathrm{Ob}\,\Phi(C)$,
$\xi = (I, (L_i)_{i\in I})$ (les $L_i \in \mathrm{Ob}\,C$), et où $\gamma$
est formé :

\begin{enumerate}
\item d'un ensemble \uncertain{$\gamma'$} (qui peut être vide) de parties
de $I$, qui sont chacune de cardinal 2 et deux à deux disjointes ;
\item pour $\forall\, J \in \gamma'$, d'un hom de
\[
\pi(\xi|J) \longrightarrow \mathbf{1}_C
\]
(où $\pi(\xi|J) = L_{j'} \otimes L_{j''}$ si $J = \{j', j''\}$)
[qui soit un accouplement parfait, i.e. une dualité].
\end{enumerate}

\note{en marge gauche, en oblique, une définition en partie illisible :
« $U^{\gamma'} = \bigcup$… »}

Les homs (\uncertain{quitte à des isoms}) de $\Phi^c(C)$ clairs. On trouve
\uncertain{des foncteurs}
\[
\mathrm{inc},\ \mathrm{res} : \Phi^{c}(C) \longrightarrow \Phi(C)
\]
(oubli des indices contractés). $\Phi^{c}(C)$ est aussi une
$\otimes$-catégorie ass. comm. unitaire ; et res est un $\otimes$-foncteur
AUC. Ceci dit, on trouve un hom canonique
\[
\Pi \circ \mathrm{inc} \xrightarrow{\ c\ } \Pi \circ \mathrm{res}
\]
(contraction) de $\otimes$-foncteurs AUC de $\Phi^{c}(C)$ dans $C$.

\page{12}
donc il n'y a pas d'isom canonique entre $L \otimes L$ et le produit de la
famille \uncertain{constante sur $I$ de valeur $L$} : il faut, pour choisir
un tel isom, choisir un ordre total sur $I$. \uncertain{C'est que}, si $I$
est de cardinal $\geq 2$, il y a a priori \ill{} possibles entre
$\pi((L_i)_{i\in I})$ et $L^{\otimes 2}$. \note{l'exposant se lit « 2 » là
où l'on attendrait « $n$ » ; le milieu de la phrase est raturé}

Mais si $L$ est inversible, la symétrie de $L \otimes L$ s'identifie à un
élément $\varepsilon(L)$ du groupe $\mathrm{Aut}(\mathbf{1}_C)$, d'ordre 2,
et le passage d'un isom. de $\Pi L$ avec $L^{\otimes n}$ à un autre, via
une permutation $\sigma$ de $[1, n]$, se fait par
$\varepsilon(L)^{\mathrm{sign}(\sigma)}$. Il n'y a aucune ambiguïté pour le
choix dans le seul cas des catégories de Picard [tout objet
\uncertain{inversible}, les $\varepsilon(L) = 1$].

Inverse d'un élément $L$ inversible : c'est un $L'$ muni d'un isom
\[
L \otimes L' \simeq \mathbf{1},
\]
\uncertain{d'où aussi}, en appliquant la symétrie, un autre isom
\[
L' \otimes L \simeq \mathbf{1},
\]
donc les relations entre $L$ et $L'$ sont symétriques ; on convient de dire
que $(L, L')$ (\uncertain{sans ordre fixé}) est un couple d'inverses si on
se donne un isom
\[
\Pi(L, L') \simeq \mathbf{1}.
\]
\note{en marge gauche, une note en oblique, en partie illisible}

\begin{tikzcd}[column sep=small, row sep=small]
L \otimes L' \arrow[r, "\varphi_{L,L'}"] \arrow[d, "s_{L,L'}"'] & \mathbf{1} \\
L' \otimes L \arrow[ur, "\varphi_{L',L}"'] &
\end{tikzcd}

\[
\varphi_{L,L'} \otimes \mathrm{id}_L \simeq \mathrm{id}_L \otimes \psi_{L',L}
\ :\ L \otimes L' \otimes L \xrightarrow{\ \sim\ } L,
\qquad \psi_{L',L} : L' \otimes L \simeq \mathbf{1}.
\]

Relations entre $\varphi_{L',L}$ et $\psi_{L',L}$ ?

Comme
\[
\varphi_{L,L'} \otimes \mathrm{id}_L\ :\ L \otimes L' \otimes L \to L,
\]
\[
\mathrm{id}_L \otimes \varphi_{L',L}
= \mathrm{id}_L \otimes (\varphi_{L,L'} \circ s_{L',L})
\ :\ L \otimes L' \otimes L \to L.
\]

Pour $q : M \to \mathbf{1}$ :

\begin{tikzcd}[column sep=small, row sep=small]
L \otimes M \arrow[r, "s_{L,M}"] \arrow[d, "L \otimes q"'] &
M \otimes L \arrow[d, "q \otimes L"] \\
L \otimes \mathbf{1} \arrow[r, "s_{L,\mathbf{1}}"'] & \mathbf{1} \otimes L
\end{tikzcd}

\note{sous le carré, les deux coins inférieurs sont reliés à $L$ par des
isoms ; à droite, trois colonnes de symboles $L\ L'\ L$ esquissent, à la
manière de ficelles, les deux contractions de $L \otimes L' \otimes L$ et
la permutation qui les échange, au facteur $\varepsilon$ près}

\section*{Théories et constructions (page 13)}

\page{13}
\note{page numérotée « 2 » de la main de l'auteur — la page 1 n'est pas
dans la chemise ; la page entière est barrée d'un long trait oblique}

On a donc
\[
\underline{\mathrm{Hom}}(I; T, T') \longrightarrow
\underline{\mathrm{Hom}}(T, T')
\]
\[
\Big\downarrow \qquad \Big\downarrow
\]
\[
\Bigl[\ \underline{\mathrm{Hom}}_0(I; T, T') \longrightarrow
\underline{\mathrm{Hom}}_0(T, T')\ \Bigr] \quad (\text{resp. } \varprojlim)
\]
\note{à droite du crochet, une précision en petit, rognée par la marge :
« où $T$, $T'$ sont dé\ldots{} »}

Les théories \uncertain{à lims} forment une 2-catégorie, soit $\mathcal{T}$
(resp. $\mathcal{T}_0$), et quand $I$ est fixé, via les
\uncertain{$I$-morphismes de théories}, une 2-catégorie $\mathcal{T}^{(I)}$
(resp. $\mathcal{T}_0^{(I)}$), et un diagramme de foncteurs de 2-catégories

\begin{tikzcd}[column sep=small, row sep=small]
\mathcal{T} \arrow[r] \arrow[d] & \mathcal{T}^{(I)} \arrow[d] \\
\mathcal{T}_0 \arrow[r] & \mathcal{T}_0^{(I)}
\end{tikzcd}

\note{la ligne qui suivait — « Enfin, si on considère les théories qui sont
définies par… », avec l'interligne « (resp. constructions infinies) » — est
biffée}

\textbf{Constructions} pour une \struck{cat} théorie $T$ : c'est par
définition un hom dans la théorie $\dot{\imath}$ correspondant à un ens. de
base sans structure, donc $\dot{\imath}_C = C$ ; donc c'est un. foncteur de
la cat. 2-fibrée sur (Cat lim finies) (resp. (Cat lim)) :
\[
T \longrightarrow \dot{\imath}.
\]

Ces constructions forment donc une catégorie
$\Gamma^{T} = \underline{\mathrm{Hom}}(T, \dot{\imath})$ (resp.
$\Gamma_0^{T} = \underline{\mathrm{Hom}}_0(T, \dot{\imath})$), qui a des
lims finies (resp. des lims) et qui est munie de
\[
\Gamma^{T} \times T \longrightarrow \dot{\imath}
\qquad (\text{resp. } \Gamma_0^{T} \times T \to \dot{\imath}_0).
\]

$C \in$ (Cat lim finies) (resp. $\in$ Cat lim) et, pour
$\forall\, \xi \in T(C)$, le foncteur
\[
u \mapsto u(\xi), \qquad \Gamma^{T} \longrightarrow C
\quad (\text{resp. } \Gamma_0^{T} \to C),
\]
commutant aux lims \uncertain{précédentes} (resp. quelc.).

\section*{Prolongement aux préfaisceaux et exactitude (page 15)}

\page{15}
\textbf{Question.} $C$ petite catégorie.
$\varphi : C \to \mathrm{Ens}$ foncteur exact à droite [$C$ stable par lim
finies]. Alors $\overline{\varphi} : \widehat{C}^{\circ} \to \mathrm{Ens}$,
qui prolonge $\varphi$ et commute aux petites lims (donné aussi par
$\overline{\varphi}(\psi) = \mathrm{Hom}(\psi, \varphi)$,
$\psi \in \mathrm{Ob}\,\underline{\mathrm{Hom}}(C, \mathrm{Ens})$),
est-il exact à droite, i.e. $\psi \mapsto \mathrm{Hom}(\psi, \varphi)$
transforme-t-il lims projectives finies en lims inductives ?

En marge gauche : $\overline{C} = \widehat{C}^{\circ} =
\underline{\mathrm{Hom}}(C, \mathrm{Ens})^{\circ}$, et le triangle

\begin{tikzcd}[column sep=small, row sep=small]
\overline{C} \arrow[dr, "\overline{\varphi}"] & \\
C \arrow[u] \arrow[r, "\varphi"'] & \mathrm{Ens}
\end{tikzcd}

\textbf{NB} \uncertain{Par l'hyp. sur $\varphi$}, c'est le cas pour les
lims finies qui proviennent de $C$. [\textbf{NB} l'inclusion
$C \to \overline{C}$ commute aux lims — finies ou non —.]

\note{une ligne biffée introduisait le cas « $C$ \struck{additive}
abélienne » ; elle s'interrompt sur « pour $\varphi = $ », et l'application
$\psi \mapsto \mathrm{Hom}(\psi, \varphi)$, « est-il exact à g. »,
$\underline{\mathrm{Hom}}(C, \mathrm{Ens}) \to (\mathrm{Ens})$ — passage
raturé, lecture d'ensemble incertaine}

\struck{$\mathrm{Hom}(\psi \times \psi_0, \varphi)$} Il faudrait que
$\varphi$ transforme mono \struck{en} \add{en} épi, i.e. que $\varphi$ soit
« \uncertain{injectif} » :

\begin{tikzcd}[column sep=small, row sep=small]
\psi \arrow[r] \arrow[dr] & \psi' \arrow[d, "\exists"] \\
& \varphi
\end{tikzcd}

Cela indique (\uncertain{prenons $\psi = \varphi$,
$\psi' = \varphi \sqcup \rho$}) que pour $\forall\, \rho : C \to
\mathrm{Ens}$, on a $\underline{\mathrm{Hom}}(\rho, \varphi) \neq
\emptyset$.

En marge gauche, le triangle correspondant :

\begin{tikzcd}[column sep=small, row sep=small]
\varphi \arrow[r] \arrow[dr] & \varphi \sqcup \psi \arrow[d, "\exists"] \\
& \varphi
\end{tikzcd}

\struck{C'est vrai si} Mais c'est \emph{toujours} faux \struck{sauf si $C$
\uncertain{vide}} : car si on prend pour $\rho$ un foncteur constant de
valeur $x$ non vide, on a $\mathrm{Hom}(\rho, \varphi) = \emptyset$, car
\[
\mathrm{Hom}_{\mathrm{Ens}}(\rho(c), \varphi(c))
= \mathrm{Hom}_{\mathrm{Ens}}(x, \emptyset) = \emptyset\ !
\]
\note{la lecture des arguments des deux membres est incertaine — le
manuscrit porte quelque chose comme $\rho(\psi_c)$, $\varphi(\psi_c)$}

\textbf{Question.} Soient $C$ petite catégorie avec lims finies, $E$
catégorie avec lims finies et petites lims inductives,
$\varphi : C \to E$ un foncteur exact à gauche,
$\overline{\varphi} : \widehat{C} \to E$ son prolongement canonique
(commutant aux petites lims inductives). Sous quelles conditions
d'exactitude sur $E$ peut-on dire que $\overline{\varphi}$ est aussi exact
à g ? [\textbf{NB} si $E = (\mathrm{Ens})^{\circ}$,
\uncertain{ça n'ira jamais, je crois} !]

On sait que c'est OK si $E$ est un topos. Réciproque ?

Supposons que $\varphi$ soit strictement \uncertain{générateur},
\struck{Alors il faut que} alors $E \to \widehat{C}$ \add{est} pl. fidèle
et $\overline{\varphi} : \widehat{C} \to E$ en est un adjoint à g.
\struck{Il suffit que $E$ soit un topos.} Donc si $\overline{\varphi}$ est
exact à g, $E$ est un topos.

\section*{Prolongement le long d'une sous-catégorie génératrice (pages 17 et 19)}

\page{17}
En tête, le dispositif : $C \subset E$, « sous-catégorie pleine
[strictement génératrice] » de « catégorie [avec petites limites
inductives] » ; $F$ « catégorie avec petites limites inductives ».
\[
\underline{\mathrm{Hom}}(C, F)
\underset{\beta}{\overset{\alpha}{\rightleftarrows}}
\underline{\mathrm{Hom}}(E, F)
\]
\[
\alpha(\varphi)(X) = \varinjlim_{Y \in C/X} \varphi(Y),
\qquad
\beta(\psi)(Y) = \psi(Y), \quad \text{i.e. } \beta(\psi) = \psi | C
= \psi \circ i.
\]
\[
\beta\alpha(\varphi)(Y) = \alpha(\varphi)(Y)
= \varinjlim_{Z \in C/Y} \varphi(Z) = \varphi(Y),
\quad \text{i.e. } \beta\alpha(\varphi) \simeq \varphi, \quad
\beta\alpha \simeq \mathrm{id}_{\underline{\mathrm{Hom}}(C,F)}.
\]
\[
\alpha\beta(\psi)(X) = \varinjlim_{Y \in C/X} \beta\psi(Y)
= \varinjlim_{Y \in C/X} \psi(Y) \longrightarrow \psi(X),
\quad \text{i.e. } \alpha\beta \to \mathrm{id}_{\underline{\mathrm{Hom}}(E,F)}.
\]
\[
\mathrm{Hom}(\alpha\varphi, \psi) \simeq \mathrm{Hom}(\varphi, \beta\psi).
\]

En fait, $\alpha = i^{F}_{!}$. Donc $i$ étant pl. fidèle, $i_{!}$ l'est.

Supposons maintenant que [$E$ \uncertain{soit aussi stable par petites
lims}], $C$ strictement génératrice, et voyons à quelle condition
$\alpha(\varphi)$ commute \uncertain{à icelles}. Comme
$\beta\alpha(\varphi) = \alpha(\varphi)|C = \varphi$, il est évidemment
nécessaire que l'on ait :

(*) Pour $\forall$ système inductif $(Y_i)$ dans $C$ dont la limite
inductive dans $E$ \struck{est} soit $Y \in \mathrm{Ob}\,C$
[\uncertain{a fortiori sa lim dans $C$ est $Y$}], on a
\[
\varphi(Y) \xleftarrow{\ \sim\ } \varinjlim \varphi(Y_i).
\]

Soit $\underline{\mathrm{Hom}}'(C, F)$ la sous-catégorie pleine de
$\underline{\mathrm{Hom}}(C, F)$ formée des $\varphi$ satisfaisant (*), et
$\underline{\mathrm{Hom}}'(E, F)$ la sous-catégorie pleine de
$\underline{\mathrm{Hom}}(E, F)$ formée des $\psi$

\page{19}
commutant aux petites lims. Donc $\beta$ induit
\[
\beta' : \underline{\mathrm{Hom}}'(E, F) \longrightarrow
\underline{\mathrm{Hom}}'(C, F).
\]
On a, pour $\psi$ dans $\underline{\mathrm{Hom}}'$,
\[
\alpha\beta(\psi)(X) = \varinjlim_{Y \in C/X} \psi(Y)
\xrightarrow{\ \sim\ } \psi(X)
\quad \text{car } X = \varinjlim_{Y \in C/X} Y,\ C
\]
étant str. gén. Donc on trouve que
\[
\beta' : \underline{\mathrm{Hom}}'(E, F) \longrightarrow
\underline{\mathrm{Hom}}'(C, F)
\]
est pleinement fidèle. (c'est une équivalence \uncertain{ssi ess. surj.})
\uncertain{Il s'agit de voir si $\alpha$ envoie bien}
$\underline{\mathrm{Hom}}'(C, F)$ dans $\underline{\mathrm{Hom}}'(E, F)$.
Or soit $(Y_i)$ un syst. inductif dans $E$, ayant dans $E$ une limite
inductive $Y$, et pour \uncertain{chaque $i$} un système inductif
$(Y_{ij})$ dans $C$, \struck{filtré sur $I$}, tel que les syst. inductifs
($i$ fixé) des $Y_{ij}$ \uncertain{aient pour limites les $Y_i$} et que le
système inductif correspondant soit $(Y_i)$. On a donc
\[
\varinjlim_i \alpha\varphi(Y_i)
= \varinjlim_i \bigl(\varinjlim_j \varphi(Y_{ij})\bigr)
= \varinjlim_{ij} \varphi(Y_{ij}),
\]
\note{le premier terme est surchargé ; la fin de la ligne, griffonnée puis
rognée par la marge, comparait ce terme à $\varphi(Y)$. Les deux lignes qui
introduisent le système double $(Y_{ij})$ sont elles-mêmes très raturées}
et la comparaison est l'hom dont \uncertain{il s'agit de} prouver que c'est
un isom. Donc on est ramené à prouver ceci :

\textbf{Lemme.} Si $\varphi \in \underline{\mathrm{Hom}}'(C, F)$, alors
pour tout $\forall$ système inductif $(Y_i)$ dans $C$ ayant une limite
inductive $Y$ dans $E$, on a
\[
\varinjlim \varphi(Y_i) \xrightarrow{\ \sim\ } \beta\varphi(Y).
\]
\note{le manuscrit porte bien $\beta\varphi(Y)$ ; on attendrait
$\alpha\varphi(Y)$, seul défini pour $Y \in E$}

Il faut donc comparer les limites pour $\varphi$ des \add{deux} systèmes
inductifs $(Y_i)$ et $C/Y \to C$. Évidemment on a :
$I \to C/Y$, cela est-il cofinal ? C'est vrai en tout cas si pour
$\forall$ diagramme fini dans $C$, la lim dans $E$ existe et est dans $C$.
\note{un point d'interrogation isolé en marge gauche face à ces deux
dernières lignes}

\end{document}
